% !TEX root = Main.tex
%
\section{Upper bound in best of both worlds} %Setting (Theorem~\ref{th:UpBOB}) }
\label{app:proofupBOB}
%
\paragraph{Stochastic setup.}
We work in the bounded i.i.d.\,stochastic setting of
Section~\ref{s:FormuBOB}, with arm laws $\nu_h$ on $[0,1]$ and
means $\mean_h$. For the analysis, relabel the arms so that
$\mean_1>\mean_2\geq\cdots\geq\mean_{\nArms}$, ordering equal
suboptimal means arbitrarily and consistently. Thus
$0<\gap_1=\gap_2\leq\cdots\leq\gap_{\nArms}\leq1$.
This does not affect the learner's rankings or tie resolution.
Fix any deterministic admissible schedule
$\propTimeVec\in\propTimeVecSpa$ and write
$\phaseTime_r=\timeHorizon\propTime_r$, so that
\[
\phaseTime_1=\phaseTime_2=\timeHorizon
\geq\phaseTime_3\geq\cdots\geq\phaseTime_{\nArms}\geq1,
\qquad \phaseTime_{\nArms+1}=0.
\]
All comparison times are integers. The schedule may depend on the
instance and horizon, but not on the realized observations.
Coincident comparison times are allowed, and sums over empty time
intervals are zero.

For $u\in[\timeHorizon]$, the rank $\tilde{\langle k\rangle}_u$
is the learner's ranking from $\estCumulGainVector_{\cdot,u-1}$.
At the unplayed time $\timeHorizon+1$, use a decreasing-score ranking
of $\estCumulGainVector_{\cdot,\timeHorizon}$ that places the actual
recommendation $\recomArm_{\timeHorizon}$ first among tied maximizers;
order other ties arbitrarily. This is only an analysis convention:
no additional pull is made, and neither the recommendation nor its
tie rule is changed.
Intuitively, after comparison time $\phaseTime_i$, we expect arms
whose means are within $\gap_i/2$ of the best mean to occupy the
first $i-1$ ranks. For $i\in[2:\nArms]$, define
\begin{equation}\label{eq:eventbob}
\eventbob_i
\triangleq
\left\{
\begin{array}{l}
\forall u\in\Integer\text{ with }\phaseTime_i<u\leq\timeHorizon+1,
\ \forall k\in\setArms:\\
\displaystyle
\mean_1-\mean_k<\frac{\gap_i}{2}
\implies\tilde{\langle k\rangle}_u<i
\end{array}
\right\}.
\end{equation}
Set $\eventbob_{\nArms+1}=\Omega$ and abbreviate
$E_i=\bigcap_{p=i+1}^{\nArms+1}\eventbob_p$.
Since only arm $1$ satisfies the eligibility condition for
$\eventbob_2$, and $\phaseTime_2=\timeHorizon$, we have
$\eventbob_2=\{\recomArm_{\timeHorizon}=1\}$.

Let $\mathcal F_{s-1}$ be the information available immediately
before drawing $\pulledArm_s$, including past actions and observed
rewards and the current resolved ranking. These information sets
form a filtration by including any next-round tie resolution in
$\mathcal F_s$, without revealing future action-sampling coins or
current rewards before sampling. Thus the ranks and the conditional
probabilities $\learnerDist_{h,s}$ are predictable. The fresh action
draw is independent of the current rewards given the past, and the
reward observed when arm $h$ is selected has conditional law $\nu_h$.
For each arm $h$, put
\[
X_{h,s}=\estGainVector_{h,s}-\mean_h,\qquad
M_h(t)=\sum_{s=1}^t X_{h,s}
=\estCumulGainVector_{h,t}-t\mean_h,\qquad
V_h(t)=\sum_{s=1}^t\Exp[X_{h,s}^2\mid\mathcal F_{s-1}].
\]
The importance-weighted estimator and $0\leq Y\leq1$ give
\begin{align*}
\Exp[X_{h,s}\mid\mathcal F_{s-1}]&=0,\\
\Exp[X_{h,s}^2\mid\mathcal F_{s-1}]
&=\frac{\Exp_{Y\sim\nu_h}[Y^2]}{\learnerDist_{h,s}}-\mean_h^2
\leq\frac1{\learnerDist_{h,s}},
\qquad
|X_{h,s}|\leq\frac1{\learnerDist_{h,s}}
\leq\nArms\bar\log\nArms.
\end{align*}
In particular, $M_h$ is a martingale and $V_h$ is its predictable
variance process; no Bernoulli assumption is needed.

\paragraph{Concentration on good-ranking events.}
We use the following fixed-time Bernstein--Freedman inequality
\citep{Freedman75OT}: for martingale differences bounded in absolute
value by $B>0$, their sum $M_t$ and predictable variance $V_t$ satisfy,
for deterministic $z,v>0$,
\[
\Pro(M_t\geq z,\ V_t\leq v)
\leq\exp\left(-\frac{z^2}{2(v+Bz/3)}\right).
\]
The same bound holds for $-M_t$. Indeed, for $0<\lambda<3/B$,
\[
\Exp[e^{\lambda X_s}\mid\mathcal F_{s-1}]
\leq\exp\left(
\frac{\lambda^2\Exp[X_s^2\mid\mathcal F_{s-1}]}
{2(1-\lambda B/3)}\right),
\]
so $\exp\{\lambda M_t-\lambda^2 V_t/[2(1-\lambda B/3)]\}$
has expectation at most one. Markov's inequality with
$\lambda=z/(v+Bz/3)<3/B$ yields the stated bound.
If a good-ranking event $E$ implies $V_t\leq v$, then
\[
\{M_t\geq z\}\cap E\subseteq\{M_t\geq z,\ V_t\leq v\},
\]
and hence the same bound applies to its probability. The event $E$
may involve observations after time $t$ or have probability zero:
we never condition the stochastic law on $E$.

\begin{lemma}\label{l:errTwoArms}
For the fixed schedule $\propTimeVec$, let $i\in[2:\nArms]$,
$j\in[i:\nArms]$, and let $k\in\setArms$ satisfy
$\mean_1-\mean_k<\gap_i/2$. For every integer $t$ with
$\phaseTime_i\leq t\leq\timeHorizon$, playing~\Pone{},
\[
\Pro\left(
\{\estCumulGainVector_{k,t}\leq\estCumulGainVector_{j,t}\}
\cap E_i\right)
\leq2\exp\left(-\frac{\timeHorizon}
{128\complexityProOne(\propTimeVec)}\right).
\]
\end{lemma}
\begin{proof}
First, the assumptions give $k<i\leq j$ and
\[
\mean_k-\mean_j
=\gap_j-(\mean_1-\mean_k)
>\gap_j-\frac{\gap_i}{2}
\geq\frac{\gap_j}{2}>0.
\]
Splitting a comparison error between the lower tail of the better
arm and the upper tail of the worse arm, we obtain
\begin{equation}\label{eq13}
\begin{split}
\Pro\left(
\{\estCumulGainVector_{k,t}\leq\estCumulGainVector_{j,t}\}
\cap E_i\right)
\leq{}&\Pro\left(E_i\cap\{-M_k(t)>t\gap_j/4\}\right)\\
&+\Pro\left(E_i\cap\{M_j(t)>t\gap_j/4\}\right).
\end{split}
\end{equation}
Define the doubled-gap rank
\[
j^+=\max\{r\in[j:\nArms]:\gap_r\leq2\gap_j\}.
\]
The set contains $j$, so $j\leq j^+\leq\nArms$ and
$\gap_{j^+}\leq2\gap_j$. If $j^+<\nArms$, then
$\gap_j<\gap_{j^++1}/2$. When $j^+=\nArms$, the universal rank
bound $\nArms$ will suffice. Thus tied gaps and exact doubled-gap
equality are included without requiring any additional arm.

We bound both terms in Equation~\ref{eq13} using the same
virtual-phase variance calculation. Write
\[
S_r=\sum_{\ell=r}^{\nArms}\ell(\phaseTime_\ell-\phaseTime_{\ell+1}),
\]
and take $(h,r)=(k,i)$ for the better arm and $(h,r)=(j,j^+)$
for the worse arm. In either case $r\geq i$ and
$t\geq\phaseTime_i\geq\phaseTime_r$.
If $r<\nArms$, the corresponding mean-gap condition is
$\mean_1-\mean_h<\gap_{r+1}/2$: for $h=k$ this follows from
$\gap_i\leq\gap_{i+1}$, and for $h=j$ it follows from the
definition of $j^+$.
On $E_i$, in each block
$\phaseTime_{\ell+1}<s\leq\phaseTime_\ell$, $r\leq\ell\leq\nArms$,
event $\eventbob_{\ell+1}$ gives
$\tilde{\langle h\rangle}_s\leq\ell$ when $\ell<\nArms$;
the bound is automatic when $\ell=\nArms$.
After $\phaseTime_r$, event $\eventbob_{r+1}$ similarly gives rank
at most $r$, with the universal bound used when $r=\nArms$.
All events used have indices larger than $i$ and belong to $E_i$.
Thus good earlier rankings increase sampling probabilities and
reduce the conditional variance. Summing the per-round bounds over
the disjoint phase blocks and the continuation interval yields,
on $E_i$,
\begin{align*}
V_h(t)
&\leq\bar\log\nArms\left[
\sum_{\ell=r}^{\nArms}\ell(\phaseTime_\ell-\phaseTime_{\ell+1})
+r(t-\phaseTime_r)\right]\\
&=\bar\log\nArms\,[S_r+r(t-\phaseTime_r)].
\end{align*}
Empty blocks contribute zero, including when comparison times coincide.

Apply the joint-event concentration bound with
$z=t\gap_r/8$, $B=\nArms\bar\log\nArms$, and the positive
variance allowance $v=\bar\log\nArms\,[S_r+r(t-\phaseTime_r)]$.
For either sign separately, it gives
\begin{align*}
\Pro\left(E_i\cap\{\pm M_h(t)\geq t\gap_r/8\}\right)
&\leq\exp\left(
-\frac{(t\gap_r/8)^2}
{2\bar\log\nArms\,[S_r+r(t-\phaseTime_r)]
+(\nArms\bar\log\nArms)t\gap_r/12}\right)\\
&=\exp\left(
-\frac{t^2\gap_r^2}
{128\bar\log\nArms\,[S_r+r(t-\phaseTime_r)
+(\nArms/24)t\gap_r]}\right).
\end{align*}
Here the factor $2$ multiplies the entire variance allowance;
$128=8^2\cdot2$, and the coefficient $\nArms/24$ comes from $Bz/3$.

The identity
$S_r=r\phaseTime_r+\sum_{\ell=r+1}^{\nArms}\phaseTime_\ell
\geq r\phaseTime_r$ implies, for $t\geq\phaseTime_r$,
\[
S_r+r(t-\phaseTime_r)\leq\frac{t}{\phaseTime_r}S_r.
\]
Consequently,
\begin{align*}
\frac{t^2\gap_r^2}
{128\bar\log\nArms\,[S_r+r(t-\phaseTime_r)
+(\nArms/24)t\gap_r]}
&\geq
\frac{\phaseTime_r^2\gap_r^2}
{128\bar\log\nArms\,[S_r+(\nArms/24)\phaseTime_r\gap_r]}\\
&\geq\frac{\timeHorizon}{128\complexityProOne(\propTimeVec)}.
\end{align*}
For the last step, the existing complexity definition, after
substituting $\phaseTime_r=\timeHorizon\propTime_r$, reads
\[
\complexityProOne(\propTimeVec)
=\timeHorizon\bar\log\nArms\,
\max_{r\in[\nArms]}
\frac{S_r+(\nArms/24)\phaseTime_r\gap_r}
{\phaseTime_r^2\gap_r^2}.
\]
Finally, $t\gap_j/4\geq t\gap_i/8$ for the better-arm tail,
and $t\gap_j/4\geq t\gap_{j^+}/8$ for the worse-arm tail.
The two one-sided joint bounds therefore control the two terms
in Equation~\ref{eq13}; summing them proves the claim for the
fixed schedule $\propTimeVec$.
\end{proof}
\begin{lemma}\label{l:errPOGerrPO}
For the fixed admissible schedule $\propTimeVec$, playing~\Pone{}
in the stochastic setting above, we have for every $i\in[2:\nArms]$,
\[
\Pro(\eventbob_i^c\cap E_i)
\leq2\nArms^2\timeHorizon
\exp\left(-\frac{\timeHorizon}
{128\complexityProOne(\propTimeVec)}\right).
\]
\end{lemma}
\begin{proof}
If $\eventbob_i$ fails, there are an integer $u$ with
$\phaseTime_i<u\leq\timeHorizon+1$ and an arm $k$ satisfying
$\mean_1-\mean_k<\gap_i/2$ but $\tilde{\langle k\rangle}_u\geq i$.
Necessarily $k<i$: otherwise $k\geq i\geq2$ and
$\mean_1-\mean_k=\gap_k\geq\gap_i$, contradicting eligibility.
At least $i-1$ arms precede $k$ in the estimated ranking, and
only $i-2$ other arms can have true labels in $[1:i-1]$.
Thus some arm $j\in[i:\nArms]$ precedes $k$. Since the ranking
uses the estimates through round $u-1$, this implies
$\estCumulGainVector_{k,u-1}\leq\estCumulGainVector_{j,u-1}$,
including when the estimated scores tie.

Since $\phaseTime_i\leq u-1\leq\timeHorizon$,
Lemma~\ref{l:errTwoArms} at $t=u-1$ gives
\[
\Pro\left(
\{\estCumulGainVector_{k,u-1}\leq\estCumulGainVector_{j,u-1}\}
\cap E_i\right)
\leq2\exp\left(-\frac{\timeHorizon}
{128\complexityProOne(\propTimeVec)}\right).
\]
There are at most $i-1$ eligible arms $k$, $\nArms-i+1$
competitors $j$, and $\timeHorizon-\phaseTime_i+1\leq\timeHorizon$
ranking times $u$, since $\phaseTime_i\geq1$.
A union bound gives
\begin{align*}
\Pro(\eventbob_i^c\cap E_i)
&\leq2(i-1)(\nArms-i+1)(\timeHorizon-\phaseTime_i+1)
\exp\left(-\frac{\timeHorizon}
{128\complexityProOne(\propTimeVec)}\right)\\
&\leq2\nArms^2\timeHorizon
\exp\left(-\frac{\timeHorizon}
{128\complexityProOne(\propTimeVec)}\right).
\end{align*}
\end{proof}%
%
	\setcounter{scratchcounter}{\value{theorem}}\setcounter{theorem}{\the\numexpr\getrefnumber{th:UpBOB}-1}
\lightres*
\setcounter{theorem}{\the\numexpr\value{scratchcounter}}
\begin{proof} %[Theorem~\ref{th:UpBOB}] 
We consider two cases separately.
	%Let us consider the $a^*_i$ that maximizes the proof quantity
	%
	%
	\paragraph{The i.i.d.~stochastic case}
For the fixed schedule $\propTimeVec$, the terminal-ranking
convention gives
\[
\ProErr_{\stoPro}(\timeHorizon)
=\Pro(\recomArm_{\timeHorizon}\neq1)
=\Pro(\eventbob_2^c).
\]
If $\eventbob_2$ fails, take the largest failed index $i\in[2:\nArms]$.
Every event with a larger index holds, since
$\eventbob_{\nArms+1}$ is certain.
Hence
\[
\eventbob_2^c\subseteq
\bigcup_{i=2}^{\nArms}(\eventbob_i^c\cap E_i).
\]
A union bound and Lemma~\ref{l:errPOGerrPO} give
\begin{align*}
\ProErr_{\stoPro}(\timeHorizon)
&\leq\sum_{i=2}^{\nArms}\Pro(\eventbob_i^c\cap E_i)\\
&\leq2(\nArms-1)\nArms^2\timeHorizon
\exp\left(-\frac{\timeHorizon}
{128\complexityProOne(\propTimeVec)}\right)\\
&\leq2\nArms^3\timeHorizon
\exp\left(-\frac{\timeHorizon}
{128\complexityProOne(\propTimeVec)}\right).
\end{align*}
This holds for every admissible schedule; its virtual comparison
times do not affect \Pone{} or the final error probability. The admissible set is finite and nonempty:
its entries lie in $\{1/\timeHorizon,\ldots,1\}$ and the all-ones
schedule is allowed. Positive comparison times and gaps make
$\complexityProOne(\propTimeVec)$ finite and positive.
Choosing a minimizing schedule gives
\[
\ProErr_{\stoPro}(\timeHorizon)
\leq2\nArms^3\timeHorizon
\exp\left(-\frac{\timeHorizon}{128\complexityProOne}\right).
\]

	\paragraph{The adversarial case}
Fix a deterministic oblivious table
$\gainVector\in[0,1]^{\nArms\times\timeHorizon}$ and relabel its
unique best arm as $1$ for analysis. Only the learner is random.
It suffices to use the sampling probability guaranteed to every arm;
no favorable ranking event is needed. With the information before
each action defined as above, the adaptive probabilities satisfy
\[
\learnerDist_{k,t}
=\Pro(\pulledArm_t=k\mid\mathcal F_{t-1})
=\frac1{\tilde{\langle k\rangle}_t\bar\log\nArms}
\geq\frac1{\nArms\bar\log\nArms}
\]
on every history. Put
\[
\cumulGain_{k,t}=\sum_{s=1}^t\gainVector_{k,s},\qquad
Y_{k,t}=\estGainVector_{k,t}-\gainVector_{k,t},\qquad
M_k^{\gainVector}(t)=\sum_{s=1}^tY_{k,s}
=\estCumulGainVector_{k,t}-\cumulGain_{k,t}.
\]
Since the table is deterministic, centering at $\gainVector_{k,t}$
is appropriate. Conditionally, the sampling indicator is Bernoulli
with parameter $\learnerDist_{k,t}$, so
\begin{align*}
\Exp[Y_{k,t}\mid\mathcal F_{t-1}]&=0,\\
\Exp[Y_{k,t}^2\mid\mathcal F_{t-1}]
&=\gainVector_{k,t}^2\left(\frac1{\learnerDist_{k,t}}-1\right)
\leq\frac1{\learnerDist_{k,t}}\leq\nArms\bar\log\nArms,\\
|Y_{k,t}|&\leq\frac1{\learnerDist_{k,t}}\leq\nArms\bar\log\nArms.
\end{align*}
Thus $M_k^{\gainVector}$ is a martingale with predictable variance
\[
V_k^{\gainVector}(\timeHorizon)
=\sum_{t=1}^{\timeHorizon}\Exp[Y_{k,t}^2\mid\mathcal F_{t-1}]
\leq\timeHorizon\nArms\bar\log\nArms
\quad\text{almost surely}.
\]

Recall that
$\timeHorizon\gap_j^{\gainVector}
=\cumulGain_{1,\timeHorizon}-\cumulGain_{j,\timeHorizon}$ for $j\neq1$
and $\gap_1^{\gainVector}=\min_{j\neq1}\gap_j^{\gainVector}$.
Hence $0<\gap_1^{\gainVector}\leq\gap_j^{\gainVector}\leq1$ and
$\complexityUnif(\gainVector)=\nArms/(\gap_1^{\gainVector})^2$.
The error event satisfies
\[
\{\recomArm_{\timeHorizon}\neq1\}
\subseteq
\{M_1^{\gainVector}(\timeHorizon)\leq-\timeHorizon\gap_1^{\gainVector}/2\}
\ \cup\!
\bigcup_{j=2}^{\nArms}
\{M_j^{\gainVector}(\timeHorizon)\geq\timeHorizon\gap_j^{\gainVector}/2\}.
\]
Indeed, outside the right-hand side, for every $j\neq1$,
\begin{align*}
\estCumulGainVector_{1,\timeHorizon}
&>\cumulGain_{1,\timeHorizon}-\timeHorizon\gap_1^{\gainVector}/2
\geq\cumulGain_{1,\timeHorizon}-\timeHorizon\gap_j^{\gainVector}/2\\
&=\cumulGain_{j,\timeHorizon}+\timeHorizon\gap_j^{\gainVector}/2
>\estCumulGainVector_{j,\timeHorizon}.
\end{align*}
Thus arm $1$ is the unique estimated maximizer, regardless of tie rules.

Apply the Bernstein--Freedman inequality above with
$B=\nArms\bar\log\nArms$, $v=\timeHorizon\nArms\bar\log\nArms$,
and $z=\timeHorizon\gap_k^{\gainVector}/2$.
The variance bound holds almost surely, so for either sign separately,
\begin{align*}
\Pro\left(\pm M_k^{\gainVector}(\timeHorizon)
\geq\timeHorizon\gap_k^{\gainVector}/2\right)
&\leq\exp\left(
-\frac{(\timeHorizon\gap_k^{\gainVector}/2)^2}
{2[\timeHorizon\nArms\bar\log\nArms
+(\nArms\bar\log\nArms)\timeHorizon\gap_k^{\gainVector}/6]}\right)\\
&=\exp\left(-\frac{\timeHorizon(\gap_k^{\gainVector})^2}
{8\nArms\bar\log\nArms(1+\gap_k^{\gainVector}/6)}\right)\\
&\leq\exp\left(-\frac{3\timeHorizon(\gap_k^{\gainVector})^2}
{28\nArms\bar\log\nArms}\right),
\end{align*}
since $1+\gap_k^{\gainVector}/6\leq7/6$.
There are $\nArms$ one-sided error events: one lower tail for arm $1$
and one upper tail for each other arm. A union bound gives
\begin{align*}
\ProErr_{\advPro(\gainVector)}(\timeHorizon)
&\leq\nArms\exp\left(-\frac{3\timeHorizon(\gap_1^{\gainVector})^2}
{28\nArms\bar\log\nArms}\right)\\
&=\nArms\exp\left(-\frac{3\timeHorizon}
{28\bar\log\nArms\complexityUnif(\gainVector)}\right)\\
&\leq\nArms\exp\left(-\frac{3\timeHorizon}
{40\bar\log\nArms\complexityUnif(\gainVector)}\right),
\end{align*}
as claimed.
\end{proof}
%
\section{On the complexities of $\complexityProOne$, $\complexitySR$, and 
	$\complexityBoth$}  %Proof of  Corollary~\ref{coco}}
\label{s:moreOnRem}
%\todomi{right reference to the remark in the title?}
%\todom{recheck the derivations below in detail}
%
%
 %\begin{remark} 
 We now show that for $\timeHorizon\geq\gap_{(\nArms)}/\gap_{(1)}$,
	$\complexityProOne = \cO\left(\complexityBoth\log^2\nArms\right)$. 
	This demonstrates that \Pone{} achieves the best 
	that can be wished for in the two worlds, up to log
	 factors.
	The extra $\log\nArms$ from the $\log^2\nArms$ is not always
	present  and we report an 
	even more detailed discussion on the three different regimes of the
	gaps used in 
	Remark~\ref{rem:LOWbobCOMP}  at the end of the section.
	%\todov{its not he right number for the corrolary numbering!}	
	\setcounter{scratchcounter}{\value{theorem}}\setcounter{theorem}{\the\numexpr\getrefnumber{coco}-1}
	\cocores*
	\setcounter{theorem}{\the\numexpr\value{scratchcounter}}
\begin{proof}%[Proof of Corollary~\ref{coco}]
Recall that $\gap_{(1)}=\gap_{(2)}$, and let
$Q=\max_{1\leq j\leq\nArms}j/\gap_{(j)}$, so that
$\complexityBoth=Q/\gap_{(1)}$. Since the rewards are in $[0,1]$,
$0<\gap_{(1)}\leq\gap_{(r)}\leq1$, $Q\geq2/\gap_{(1)}$, and $Q\geq\nArms$.
Set the ideal fractions $\alpha_r=\gap_{(1)}/\gap_{(r)}$ and the rounded fractions
$\propTime_r=\lceil\timeHorizon\alpha_r\rceil/\timeHorizon$ for
$r\in\setArms$, with $\alpha_{\nArms+1}=\propTime_{\nArms+1}=0$.
Ceilings preserve weak ordering, so
$1=\propTime_1=\propTime_2\geq\cdots\geq\propTime_{\nArms}>0$ and
$1\leq\timeHorizon\propTime_r\leq\timeHorizon$, with
$\timeHorizon\propTime_r\in\Integer$.
Thus $\propTimeVec\in\propTimeVecSpa$, including when comparison rounds coincide.

For an ordered positive vector $b$ with $b_1=b_2=1$ and $b_{\nArms+1}=0$, let
\[
\Phi(b)=\bar\log\nArms\,
\max_{1\leq r\leq\nArms}
\frac{\sum_{j=r}^{\nArms}j(b_j-b_{j+1})
+(\nArms/24)b_r\gap_{(r)}}{b_r^2\gap_{(r)}^2}.
\]
This is the algebraic expression defining $\complexityProOne(\propTimeVec)$;
we use it for nonintegral comparison times only as an algebraic comparison.
The identity
\[
\sum_{j=r}^{\nArms}j(b_j-b_{j+1})
=r b_r+\sum_{j=r+1}^{\nArms}b_j
\]
makes all coefficients in the numerator nonnegative.

Since $\alpha_r\gap_{(r)}=\gap_{(1)}$, the second numerator term divided by
the denominator is $\nArms/(24\gap_{(1)})\leq\complexityBoth/24$.
For the first term, when $r\geq2$, the identity gives
\begin{align*}
\frac{1}{\gap_{(1)}}\sum_{j=r}^{\nArms}j(\alpha_j-\alpha_{j+1})
&=\frac{r}{\gap_{(r)}}+\sum_{j=r+1}^{\nArms}\frac{1}{\gap_{(j)}}\\
&\leq Q\left(1+\sum_{j=r+1}^{\nArms}\frac{1}{j}\right)
\leq Q\left(\bar\log\nArms-\frac12\right).
\end{align*}
For $r=1$, the weighted sum equals that for $r=2$ because $\alpha_1=\alpha_2$.
The tail sums are empty when $\nArms=2$, so this case is also covered.
It follows that
\[
\Phi(\alpha)\leq\complexityBoth\bar\log\nArms
\left(\bar\log\nArms-\frac12+\frac{1}{24}\right)
\leq\complexityBoth(\bar\log\nArms)^2.
\]

If $\timeHorizon\geq\gap_{(\nArms)}/\gap_{(1)}$, then
$\timeHorizon\alpha_r\geq1$ for every rank, and
\[
\alpha_r\leq\propTime_r<\alpha_r+\frac{1}{\timeHorizon}\leq2\alpha_r.
\]
Moving each comparison by less than one round therefore makes the telescoped
numerator at most twice its value at $\alpha$, while the denominator is
no smaller. Using the admissible rounded schedule, we obtain
\[
\complexityProOne\leq\Phi(\propTimeVec)\leq2\Phi(\alpha)
\leq2\complexityBoth(\bar\log\nArms)^2
=\cO\left(\complexityBoth\log^2\nArms\right),
\]
where the last comparison uses $\bar\log\nArms\leq1+\log\nArms$.
Theorem~\ref{th:UpBOB} then gives
\[
\ProErr_{\stoPro}(\timeHorizon)
\leq2\nArms^3\timeHorizon\exp\left(-\frac{\timeHorizon}{128\complexityProOne}\right)
\leq2\nArms^3\timeHorizon\exp\left(
-\frac{\timeHorizon}{256\complexityBoth(\bar\log\nArms)^2}\right).
\]
For $1\leq\timeHorizon<\gap_{(\nArms)}/\gap_{(1)}$,
$\complexityBoth\geq2/\gap_{(1)}^2$ gives
\[
\frac{\timeHorizon}{256\complexityBoth(\bar\log\nArms)^2}
<\frac{\gap_{(1)}\gap_{(\nArms)}}{512(\bar\log\nArms)^2}\leq\frac{1}{512}.
\]
Since $2\nArms^3\timeHorizon\geq16$ and $\exp(-x)\geq1-x$, the displayed
error bound exceeds one and follows from $\ProErr_{\stoPro}(\timeHorizon)\leq1$.
This proves the error bound for all positive integer horizons, without
asserting the complexity comparison for these smaller budgets.
\end{proof}
%

\subsection{Relation between 
	$\complexityProOne$, $\complexitySR$, and 
	$\complexityBoth$ for different regimes of the 
	gaps.}
We now study the relation between 
$\complexityProOne$, $\complexitySR$, and 
$\complexityBoth$ for different regimes of the 
gaps.
We use the same examples as the ones used in 
Remark~\ref{rem:LOWbobCOMP}.
We first choose ideal fractions $b$,
with $1=b_1=b_2\geq\cdots\geq b_{\nArms}>b_{\nArms+1}=0$,
and use the algebraic expression $\Phi$ from the preceding corollary proof.
Set $\propTime_r=\lceil\timeHorizon b_r\rceil/\timeHorizon$ for $r\in\setArms$
and $\propTime_{\nArms+1}=0$. Ceilings preserve weak ordering and give
admissible integer times, even when some coincide; only the analysis
changes, while \Pone{} and its horizon are unchanged.
If $\timeHorizon b_{\nArms}\geq1$, the positive telescoped numerator gives
$\complexityProOne\leq\Phi(\propTimeVec)\leq2\Phi(b)$, as proved above.
These schedules compare larger-gap arms earlier, keeping the best two
arms at the final comparison time.
In the stated horizon ranges, the comparisons are at worst
$\cO(\complexityBoth\log^2\nArms)$, and
$\cO(\complexityBoth\log\nArms)$ in the flat and square-root regimes.

\paragraph{\raisebox{.04cm}{\textcolor{bull}{$\blacktriangleright$}}~Flat regime}  All the gaps are equal, that is,
$k\in\setArmsmo,$ we have that $\gap_k=\gap_{1}$. 

\smallskip
\noindent
We choose $b_r=1$, $\forall r\in\setArms$.
All comparisons occur at the final round, and no rounding is needed.
The weighted sum equals $\nArms$ at every rank, and
$\complexitySR=\complexityBoth=\nArms/\gap_{(1)}^2$. Thus the two terms give
\[
\complexityProOne\leq\Phi(b)
=\bar\log\nArms\left(\frac{\nArms}{\gap_{(1)}^2}+\frac{\nArms}{24\gap_{(1)}}\right)
\leq\frac{25}{24}\complexityBoth\bar\log\nArms,
\]
since $\gap_{(1)}\leq1$. The single-logarithmic comparison holds for every
positive integer horizon.

\paragraph{\raisebox{.04cm}{\textcolor{bull}{$\blacktriangleright$}}~Super-linear gaps}  Since $2\in\argmin_{r\in\setArms}(\gap_{(r)}/r)$,
$\gap_{(r)}\geq r\gap_{(1)}/2$ for $r\geq2$ and
$\complexitySR=\complexityBoth=2/\gap_{(1)}^2$.
Also, $\nArms\gap_{(1)}\leq2$ since $\gap_{(\nArms)}\leq1$.

\smallskip
\noindent
Let $b_1=1$ and $b_r=2/r$, $r\in[2:\nArms]$.
This normalizes the reciprocal-rank schedule, with $b_2=1$ and
$b_r\gap_{(r)}\geq\gap_{(1)}$. The large gaps permit early identification
of bad arms.
The second complexity term satisfies
\[
\frac{\nArms}{24b_r\gap_{(r)}}
\leq\frac{\nArms}{24\gap_{(1)}}\leq\frac{\complexityBoth}{24}.
\]
For the first term, telescoping gives, for $r\geq2$,
\[
\sum_{j=r}^{\nArms}j(b_j-b_{j+1})
=2+2\sum_{j=r+1}^{\nArms}\frac1j
\leq2\left(\bar\log\nArms-\frac12\right).
\]
The rank-$1$ sum equals the rank-$2$ sum since $b_1=b_2$.
Using $b_r\gap_{(r)}\geq\gap_{(1)}$ gives
\[
\Phi(b)\leq\complexityBoth\bar\log\nArms
\left(\bar\log\nArms-\frac12+\frac1{24}\right)
\leq\complexityBoth(\bar\log\nArms)^2.
\]
Since $b_{\nArms}=2/\nArms$, rounding gives
\[
\complexityProOne\leq2\complexityBoth(\bar\log\nArms)^2
=\cO(\complexityBoth\log^2\nArms),\qquad\timeHorizon\geq\nArms/2.
\]

\paragraph{\raisebox{.04cm}{\textcolor{bull}{$\blacktriangleright$}}~Square-root gaps} We have that $2\in\argmin_{r\in\setArms}(\gap^2_{(r)}/r)$,
$\sqrt{\nArms/2}\gap_{(1)}=\gap_{(\nArms)}$, and also that
$\sqrt{r/2}\gap_{(1)}\geq\gap_{(r)}$ for each intermediate rank $3\leq r<\nArms$.
The minimizing condition gives the reverse inequality, hence
$\gap_{(r)}=\gap_{(1)}\sqrt{r/2}$ for every $r\geq2$.
Thus $\complexitySR=2/\gap_{(1)}^2$ and
$\complexityBoth=\sqrt{2\nArms}/\gap_{(1)}^2=\sqrt{\nArms/2}\complexitySR$.

\smallskip
\noindent
Let $b_1=1$ and $b_r=\sqrt{2/r}$, $r\in[2:\nArms]$.
This normalizes the square-root schedule, with $b_2=1$ and
$b_r\gap_{(r)}=\gap_{(1)}$.
For $r\geq2$, the integral comparison
$\sum_{j=r+1}^{\nArms}j^{-1/2}\leq2(\sqrt{\nArms}-\sqrt r)$ gives
\[
\sum_{j=r}^{\nArms}j(b_j-b_{j+1})
=\sqrt{2r}+\sum_{j=r+1}^{\nArms}\sqrt{2/j}
\leq2\sqrt{2\nArms}.
\]
The rank-$1$ sum equals the rank-$2$ sum. The first complexity term
is at most $2\complexityBoth$, and the second is
$\nArms/(24\gap_{(1)})\leq\complexityBoth/24$, since
$\gap_{(\nArms)}=\gap_{(1)}\sqrt{\nArms/2}\leq1$.
Thus $\Phi(b)\leq(49/24)\complexityBoth\bar\log\nArms$.
Since $b_{\nArms}=\sqrt{2/\nArms}$, rounding gives
\[
\complexityProOne\leq2\Phi(b)\leq\frac{49}{12}\complexityBoth\bar\log\nArms
=\cO(\complexityBoth\log\nArms),\qquad
\timeHorizon\geq\sqrt{\nArms/2}.
\]
For $\nArms=2$, the tail sums are empty and both normalized schedules
reduce to the flat case, without rounding loss.
The larger $\complexityBoth$ reflects the extra variance incurred before
focusing on the best arms, as discussed in Remark~\ref{rem:LOWbobCOMP}.
For shorter horizons, these complexity bounds are not asserted.
In either nonflat regime the threshold is at most $1/\gap_{(1)}$ and
$\complexityBoth\geq2/\gap_{(1)}^2$, so $\timeHorizon/\complexityBoth<1/2$.
The corresponding derived probability bounds from Theorem~\ref{th:UpBOB} exceed one:
the exponent has magnitude less than $1/256$ and
$2\nArms^3\timeHorizon\geq16$. They follow from
$\ProErr_{\stoPro}(\timeHorizon)\leq1$, as in Corollary~\ref{coco}'s proof.
%\todomi{gramarly it all in the end}
